\section{The African Ass}

\begin{quotex}
Modern civilizations are eaters of space and traditional civilizations were eaters of time.
\flright{\textsc{Julius Evola}, \emph{Civilization of Space and Civilization of Time}}
\end{quotex}

\textbf{Julius Evola} explains that ancient civilizations were stable, static, and ultraconservative while modern civilizations claim to be progressive and dynamic. In effect, that means the latter are subject to contingency, to the movement of incessant change, and to a rapid ascent and an equally rapid vertiginous decline. The ancient rishis, or seers, had a clairvoyant vision of the supernatural, the eternal, the immutable. But the new visionaries are shown in motion, with no particular destination.

Thus the first visionary was in a taxicab in random motion, and the next visionary professor is shown walking through an airport. We are given to understand he is traveling, but we don't know where, just that he is visiting the world. His father was British and his mother Ghanaian. When asked about that unlikely pairing, his father's simple answer was, ``I liked her African ass, it was so different from the English ass.'' Is there any wonder that his son became such a deep thinker?

While, by tradition, he takes on the race of his father, Evola points out such crossings can cause an imbalance in the life of the soul. Sure enough, this was the theme of his vision. Since he was not confined to a village, he regarded himself as a ``cosmopolitan'', that is, nothing in particular. We all recall \textbf{Joseph de Maistre}'s objection to this:

\begin{quotex}
But there is no such thing as man in the world. In my lifetime I have seen Frenchmen, Italians, Russians, etc.; thanks to Montesquieu, I even know that one can be Persian. But as for man, I declare that I have never in my life met him; if he exists, he is unknown to me.
\end{quotex}


Maistre is thinking in terms of space, but for a man in time, such labels are apparently meaningless.

Yet, reality impinges even on this cosmopolitan vision, since, in his travels among the less privileged, he encounters the particular. Hence, his conundrum: how to balance cosmopolitanism with diversity? As an example, he points out the matriarchal structure of the Ghanaian village. In such a system, the male caretaker of children is not the father, but rather the cross-uncle (the mother's brother). Genetically, it makes sense if you have doubts about your wife's fidelity. Psychologically, it is more difficult to think of your mother that way; even so, you are assured of some genetic sharing with your sister, and hence her children.

But how would such diversity work in practice. For example, in the UK the father is held responsible and can be sued for support. What if the father's defense is that in his culture that is not so, and they should be going after his brother-in-law. Do you see, then, the dilemma of the cosmopolitan? Nevertheless, the professor believes he can have both. The professor suffers from the double delusion of the modern mind. First, there is ignorance of metaphysical and traditional principles. Second, there is the arrogance to presume he can dispense with them altogether. Since reality is unavoidable, the professor has stumbled onto something, but it is better understood in the light of \textbf{Empedocles}.

First of all, the cosmos in cosmopolitan does not refer to the world in toto, but rather to the ordered world. Thus when John writes, ``God so loved the cosmos'', he tells us that God loves the world that is ordered through the Logos and subject to the Law. This is the love, the force of attraction that binds a community together. Strife, then, is the force of dispersion that unbinds the community. Diversity, therefore, needs to be integrated into that cosmic order; only with this understanding can one be both a cosmopolitan and a promoter of diversity. Contrarily, to conceive it on egalitarian grounds is to make it part of Strife. This is not a moral judgment but a metaphysical deduction. For those committed to continuity of space, disordered diversity represents a destruction of their tradition. For the modern mind, lost in time, celebrating instinct, novelty, sensation, and materiality, this is the mirage of some future greatness.


\flrightit{Posted on 2011-09-07 by Cologero}

\begin{center}* * *\end{center}

\begin{footnotesize}
\begin{sffamily}
\texttt{logres on 2011-09-08 at 08:13 said:}

Rosenstock Huessy said that man's great desire was to segregate space and unite time, whereas the modern world wishes to do the exact opposite. Thank you. This explains exactly his insight, in a traditional framework. It also explains the ``Gospel'' in a manner I have never heard -- God can't love chaos.

\hfill

\texttt{Will on 2011-09-09 at 08:50 said:}

Readers may be interested to know that the ``philosophers'' referred to in this post and in The Crazy Uncle are from the film Examined Life, which is comprised of interviews with contemporary theorists, and which perfectly illustrates the vast chasm between ancient philosophy and its modern counterfeit.

\hfill

\texttt{Cologero on 2011-09-09 at 12:19 said:}

Thanks for digging that up. I'm surprised you were able to recognize it, given my rather tendentious description of the vignettes.


\hfill

\end{sffamily}
\end{footnotesize}

